\origpage{332}{N. Robertson, P.D. Seymour / Journal of Combinatorial Theory, Series B 90 (2004) 325--385}
A bad sequence for a colour scheme $\chi$ is a countable sequence $(\Gamma_i,\phi_i)$ $(i=1,2,\dots)$ of $\chi$-coloured paintings, such that for all $j>i\ge 1$ there is no linear inflation of $(\Gamma_i,\phi_i)$ in $(\Gamma_j,\phi_j)$. We say that a colour scheme is bad if there is a bad sequence for it. Thus, we wish to show that no colour scheme is bad. We wish to consider the following four statements $\mathbf{S}_1\text{--}\mathbf{S}_4$ about a colour scheme $\chi$.

A surface without boundary is \textit{simpler} than another if the second can be obtained from the first by adding at least one handle or crosscap. Our first statement is

$\mathbf{S}_1$ \textit{There is no bad colour scheme $\chi'$ with $\hat{\Sigma}_{\chi'}$ simpler than $\hat{\Sigma}_\chi$.}

A set $\{(\Gamma_i,\phi_i):i\in I\}$ of $\chi$-coloured paintings is said to be \textit{similarly oriented} if either $\Sigma_\chi$ is not orientable, or there is an orientation $\omega$ of $\Sigma_\chi$ such that for all $i\in I$ and every $e\in E(\Gamma_i)$ with $|\tilde{e}|=3$, the orientations of $e$ given by $\omega$ and given by the cyclic order of the terms of $\gamma_{\Gamma_i}(e)$ coincide. A colour scheme $\chi$ is \textit{orientedly bad} if there is a similarly oriented bad sequence for it.

If $\Omega,\Omega'$ are quasi-orders we write $\Omega'\preceq \Omega$ if

\hspace*{12pt}$\bullet$ $E(\Omega')\subseteq E(\Omega)$,

\hspace*{12pt}$\bullet$ for each $x'\in E(\Omega')$ and $x\in E(\Omega)$, if $x\le x'$ in $\Omega$ then $x\in E(\Omega')$, and

\hspace*{12pt}$\bullet$ for $x_1,x_2\in E(\Omega')$, $x_1\le x_2$ in $\Omega'$ if and only if $x_1\le x_2$ in $\Omega$.

We write $\Omega'\prec \Omega$ if $\Omega'\preceq \Omega$ and $\Omega'\ne \Omega$. We use $\cong$ to denote homeomorphism. Our second statement is

$\mathbf{S}_2$ \textit{There is no orientedly bad colour scheme $\chi'$ with $\hat{\Sigma}_{\chi'}\cong \hat{\Sigma}_\chi$ such that $\Omega_{\chi'}(3)\preceq \Omega_\chi(3)$ and $\Omega_{\chi'}(2)\preceq \Omega_\chi(2)$ where at least one of the inclusions is strict.}

We denote the number of cuffs of a surface $\Sigma$ by $c(\Sigma)$.

$\mathbf{S}_3$ \textit{There is no orientedly bad colour scheme $\chi'$ with $\hat{\Sigma}_{\chi'}\cong \hat{\Sigma}_\chi$, $\Omega_{\chi'}(k) = \Omega_\chi(k)$ $(k = 2,3)$, and $c(\Sigma_{\chi'}) < c(\Sigma_\chi)$.}

Now let $\chi,\chi'$ be colour schemes. We say that $\chi'$ is a \textit{refinement} of $\chi$ if $\Sigma_{\chi'}\cong \Sigma_\chi$, $\Omega_{\chi'}(k) = \Omega_\chi(k)$ $(k = 2,3)$, and there is a function $f$ from the set of long sides of $\Phi_{\chi'}$ to the set of long sides of $\Phi_\chi$, such that

\hspace*{12pt}$\bullet$ $\Omega_{\chi'}(R)\preceq \Omega_\chi(f(R))$ for each long side $R$ of $\Phi_{\chi'}$,

\hspace*{12pt}$\bullet$ if $R_1,R_2$ are distinct long sides of $\Phi_{\chi'}$ with $f(R_1) = f(R_2)$ then $\Omega_{\chi'}(R_1)\prec \Omega_\chi(f(R_1))$,

\hspace*{12pt}$\bullet$ if for each long side $S$ of $\Phi_\chi$ there is a long side $R$ of $\Phi_{\chi'}$ with $f(R) = S$ and $\Omega_{\chi'}(R) = \Omega_\chi(S)$, then $\Phi_{\chi'}$ has fewer short sides then $\Phi_\chi$.

In this case, we call $f$ an \textit{embedding} of $\chi'$ in $\chi$.

$\mathbf{S}_4$ \textit{There is no orientedly bad colour scheme which is a refinement of $\chi$.}
